\documentclass[11pt]{article}
\usepackage[T1]{fontenc}
\usepackage{lmodern}
\usepackage{microtype}
\usepackage[margin=1in]{geometry}
\usepackage{amsmath,amssymb,amsthm,mathtools}
\usepackage{algorithm}
\usepackage{algpseudocode}
\usepackage{booktabs}
\usepackage[numbers,sort&compress]{natbib}
\usepackage[hidelinks]{hyperref}
\usepackage[capitalise,noabbrev]{cleveref}
\setlength{\emergencystretch}{3em}

\theoremstyle{plain}
\newtheorem{theorem}{Theorem}[section]
\theoremstyle{definition}
\newtheorem{definition}[theorem]{Definition}

\newcommand{\lean}[1]{\path{#1}}
\newcommand{\N}{\mathbb{N}}
\newcommand{\at}[2]{#1|_{#2}}
\newcommand{\btw}[3]{#1|_{(#2,\,#3]}}
\newcommand{\hd}[1]{\lvert #1 \rvert}
\newcommand{\pre}{\preceq}
\newcommand{\cat}{\mathbin{+\!\!+}}

\title{The Landing Log:\\
A Catalog's Past in Its Own File, with Holes That Keep Their Keys}
\author{Alberto Serrano-Calva}
\date{November 11, 2026}

\begin{document}
\maketitle

\begin{abstract}
The catalog of a knowledge store maps each path to the content it holds now and is rewritten in
place, so it has no past: nothing records which content a path held before, or whether it was
deleted. A design for reading the store as of a past cut found this the one missing ledger, and
this paper builds it as the design's first stage. The landing log is two tables in the catalog's
own file, appended to by the catalog's six writing methods in the transaction of the row each
writes, so the catalog is the fold of the log and no caller can write one without the other. A
row carries a salted commitment that a hash chain covers, so a purge empties a row and the chain
still verifies; a row names its path by a keyed hash whose key is split between the file and a
secret outside it, so a hole still shadows the path's earlier versions, and the file alone cannot
test a guess of a purged path. The design takes two forms of removal: a retire leaves the item's
history readable at the cuts where it existed, and a purge leaves a hole at every cut. Five laws
of the design (append invariance, monotonicity, composition, idempotence, and a derived view as a
function of the prefixes) enter the store's ledger of laws as 47 statements machine-checked in
Lean~4, and a property test checks the code's reads against them over random sequences of writes.
On a copy of a catalog of 400{,}274 items the migration that writes the log takes 6.3~s and
429~bytes a row, the log adds 42 to 69~$\mu$s to a key's write, and a read of one path at a cut
takes 15~$\mu$s. Measuring at that size also found a query plan that walked the whole chain on
every write.
\end{abstract}

\section{Introduction}

The system is a personal knowledge store \cite{serranocalva2026palace}: notes, dialogue, papers and
code land through one writer into a dozen stores, and a sentence of a query algebra reads them.
Among the stores is a catalog, a table from each source path to the digest of the content it
holds, its title, its provenance class, an identity, and whether it is active. The writer reads it
to decide that content is unchanged, writes it last as the idempotence key of a landing, and
flips a row inactive when its source deletes the item. It is rewritten in place. A design for
reading the store as of a past cut \cite{serranocalva2026palace} found that this is the one store
whose past nothing else holds: the version chains record the past of one place only, and the
attestation ledger names digests and no paths. So no read at a past cut can say which version of
a note a cut holds.

This paper builds the first stage of that design, the landing log, and reports what it costs.
The contributions:
\begin{enumerate}
\item The log (\Cref{sec:log}): an append-only, hash-chained record of the catalog's changes, one
  chain a provenance class, written by the catalog's own methods in the transaction of the row
  each writes, with the reads a later stage needs (the state of a path at a cut, its history, every
  head at a cut, a range between cuts) and a verifier.
\item The rules that keep an erasure compatible with the chain and with reads by path: a purge
  empties a row and keeps its position, its commitment and its hash, and the row keeps a keyed
  hash of its path, whose key is half in the file and half in a secret outside it.
\item Five laws (\Cref{sec:laws}), proved in Lean over the design's model, now stated in the
  store's ledger and judged by a comparator, with a property test of the code against each
  (\Cref{sec:eval}).
\item Measurements on a copy of a 400{,}274-item catalog (\Cref{sec:eval}), and three places where
  building the design changed it.
\end{enumerate}

\paragraph{A decision of the design: two forms of removal.}
A retire, the tombstone a deletion in the source writes, is the shallow form: the item leaves the
present, and its history stays readable at the cuts where it existed; no seal withholds it there,
and a note written later at the same path takes a fresh identity, so that it does not inherit the
retired note's versions. A purge is the deep form: a hole at every cut, and every derived answer
at a past cut becomes what it would have been had the item never landed; what survives is the
record that a purge happened, not the content. The design's owner took this decision on 11
November 2026, after the design was reviewed. The log holds both: a retire is a row like any
other, and a purge empties rows. The fresh identity is a later stage's.

\section{Background and related work}

\paragraph{Hash chains and audit logs.}
Linking each record to the hash of its predecessor makes an append history tamper-evident
\cite{haber1991timestamp,merkle1988signature}, and secure audit logs add keys so that an attacker
who later holds the log cannot rewrite or read what came before \cite{schneier1999audit}. Here the
chain covers a salted commitment to each row and not the row \cite{halevi1996commitment}, so a row
can be emptied and the chain kept; once the salt is gone the commitment cannot be tested against a
guess of the content. A plain digest of guessable data is such a test, the reason salts were
added to stored passwords \cite{morris1979password}. The key of a path is an HMAC
\cite{bellare1996keying,rfc2104}.

\paragraph{Erasure in append-only records.}
A right to erasure meets append-only records in many systems \cite{politou2018forgetting}, and one
answer for a hash-chained ledger is to make blocks redactable \cite{ateniese2017redactable}. The
log needs less: one writer, and a chain over commitments that an erasure leaves in place.

\paragraph{Temporal data.}
Temporal databases keep rows valid over periods \cite{snodgrass1985taxonomy}, and SQL:2011's
system-versioned tables answer a query as of a system time \cite{kulkarni2012sql2011}; multiversion
concurrency control keeps versions so a reader sees a past state \cite{bernstein1983mvcc}. The log
counts periods in positions of an append chain, never in wall time, since the store's stores share
no clock; a cut is a position a chain, after the consistent cuts of distributed systems
\cite{lamport1978time,chandy1985snapshots}.

\paragraph{Checking a transcription.}
The laws are proved in Lean~4 \cite{moura2021lean4} over Mathlib \cite{mathlib2020}, on a model, and
the code is checked against them by property-based testing \cite{claessen2000quickcheck} with
Hypothesis \cite{maciver2019hypothesis}.

\section{The model and the laws}\label{sec:laws}

\begin{definition}[Stores, cuts, reads]
A \emph{store} is a family of append chains $s : K \to \mathrm{List}\,\alpha$, its \emph{head}
$\hd{s}(k)$ the length of chain $k$; $s \cat t$ appends chain by chain, and $s \pre s'$ when
every chain of $s$ is a prefix of its chain in $s'$. A \emph{cut} is $c : K \to \N$, ordered
pointwise. The \emph{read at} $c$ is $\at{s}{c}(k) = \mathrm{take}\,(c\,k)\,(s\,k)$, the
\emph{read between} $a$ and $b$ is $\btw{s}{a}{b}(k) = \mathrm{drop}\,(a\,k)\,(\at{s}{b}\,k)$, and
a \emph{derived view} at $c$ is $D(\at{s}{c})$. A reader $F(s,c)$ is \emph{append invariant} when
$F(s \cat t, c) = F(s, c)$ for every $t$ and every $c \le \hd{s}$.
\end{definition}

\begin{definition}[A keyed chain, three ways]
For a chain $l$ of keyed writes $(x, v)$: $\mathrm{stateAt}\,l\,c$ maps each key to its last value
in $\mathrm{take}\,c\,l$; $\mathrm{lastWrite}\,l\,c\,x$ is the value of the last write of $x$ in
that prefix, one seek; and $\mathrm{gens}\,l$ is one row a write, valid from the first cut that
holds it up to the first cut that holds the key's next write, its \emph{closing stamp}, read at a
cut by $\mathrm{genRead}$.
\end{definition}

\begin{theorem}[SNP-1, append invariance]\label{thm:snp1}
For $c \le \hd{s}$: $\at{(s \cat t)}{c} = \at{s}{c}$, a read between two cuts inside the store is
unchanged likewise, and every derived view is append invariant. So are the read under an erasure
and under a grant, the keyed state and the generations (whose closing stamps growth writes), the
history of a key with no seal, and the filtered read of the references a resolver does not
follow. The controls: a reader that ignores the cut is not append invariant, and a retire written
after a cut empties the history a seal serves there.
(\lean{readAt_grow}, \lean{readBetween_grow}, \lean{appendInvariant_view},
\lean{stateAt_append}, \lean{genRead_gens_append}, \lean{landingsAt_append},
\lean{kept_grown_iff}; \lean{Witness.present_not_appendInvariant},
\lean{Witness.seal_empties_history}.)
\end{theorem}

\begin{theorem}[SNP-2, monotonicity]\label{thm:snp2}
If $c \le d$ then $\at{s}{c} \pre \at{s}{d}$; $\at{(\at{s}{c})}{d} = \at{s}{(d \sqcap c)}$ and
$\hd{\at{s}{c}} = c \sqcap \hd{s}$. The history of a key with no seal and the filtered read grow
with the cut. (\lean{readAt_mono}, \lean{readAt_readAt}, \lean{head_readAt},
\lean{landingsAt_mono}, \lean{Kept.mono}.)
\end{theorem}

\begin{theorem}[SNP-3, composition]\label{thm:snp3}
If $a \le b$ then $\at{s}{a} \cat \btw{s}{a}{b} = \at{s}{b}$, and two adjacent reads-between
compose. So a fold continued from a checkpoint over the delta is the fold of the prefix, for a
store and for a keyed chain, and a key the delta does not write has one state at both cuts.
(\lean{readAt_append_readBetween}, \lean{readBetween_append}, \lean{foldAt_readBetween},
\lean{stateAt_eq_foldl}, \lean{stateAt_eq_of_not_written}.)
\end{theorem}

\begin{theorem}[SNP-4, idempotence]\label{thm:snp4}
If $\hd{s} \le c$ then $\at{s}{c} = s$; every derived view at the head is the derivation of the
store, and the generations at the head read the store rewritten in place.
(\lean{readAt_of_head_le}, \lean{readAt_head}, \lean{view_head}, \lean{genRead_gens_length}.)
\end{theorem}

\begin{theorem}[SNP-5, a derived view is a function of the prefixes]\label{thm:snp5}
A reader is append invariant if and only if it reads only the prefixes, if and only if it agrees
with a derived view on every cut inside the store. For a keyed chain,
$\mathrm{genRead}\,(\mathrm{gens}\,l)\,c\,x = \mathrm{stateAt}\,l\,c\,x =
\mathrm{lastWrite}\,l\,c\,x$, and at most one generation of a key is alive at a cut. The control:
the store rewritten in place determines no read at a past cut.
(\lean{appendInvariant_iff_prefixOnly}, \lean{appendInvariant_iff_exists_view},
\lean{genRead_gens}, \lean{stateAt_eq_lastWrite}, \lean{length_filter_selects_le_one};
\lean{Witness.stateOf_forgets}.)
\end{theorem}

The proofs are short inductions on lists; the substance is in the choice of statements. SNP-5 is
what lets the log be read three ways that must agree: by a seek on a path's key, by the closing
stamp for every path at once, and by folding the prefix. The log's \lean{closed_at} column is the
closing stamp, and it is the one field written after a row is appended; SNP-1 for the generations
says that writing it changes no read at a cut inside the chain. The design's other five laws
(the grant, the erasure, resolution, certificates, the horizon) stay design theorems until the
stage that relies on each is built; SNP-7's holes and SNP-9's commitments shape the log below, and
the tests of \Cref{sec:eval} check what of them the log already does.

\section{The log as built}\label{sec:log}

\paragraph{Tables.}
\lean{landing_chains} has one row a chain: an integer id never reused, the class it holds now, the
chain's 32 random bytes, and a key check. \lean{landings} has one row an event, keyed by
\lean{(chain, seq)}: its kind, the path's key, the path, the identity, the version number where
the place keeps version chains, the digest and title of a landing, a horizon's position, a 16-byte
salt, the commitment, the hash, the closing stamp, and a wall-clock bookmark that orders nothing.
Three partial indexes serve the reads: by path key, by identity where it differs from the path, and
by closing stamp.

\paragraph{Kinds.}
A \lean{store} is a landing; a \lean{retire} the tombstone of an identity no other active path
holds; a \lean{move} a tombstone that leaves the identity at another path (the second half of a
rename), or an identity taken off a path by an owner's re-key; a \lean{leave} the row that closes
a path in its chain when the path changes class, its \lean{store} then landing in a chain of the
new class; \lean{genesis} and \lean{horizon} open a chain. A write that changes nothing appends
nothing. A relabel of a whole class rewrites the class of its chains and no row, so no position
moves and no hash breaks.

\paragraph{The fold.}
The catalog is the fold of the log. A path's present is its last event in the one chain where that
event is not a \lean{leave}; a \lean{store} there is an active row, a \lean{retire} or a
\lean{move} a tombstoned row with the digest and title of the path's last \lean{store}, and a
\lean{store} emptied by a purge no row. The design had a \lean{move} for a change of class as
well, and read every \lean{move} as no row. The catalog contradicts that reading: it keeps a renamed
path's row, tombstoned. With one kind, a path that changed class and was then renamed has an open
\lean{move} in two chains and no way to tell which is its present; with \lean{leave} the present is
unique, also after a purge, since a purge keeps a row's kind.

\paragraph{Commitments, holes and keys.}
The commitment is SHA-256 over the salt and the canonical JSON of the row's fields; the hash is
SHA-256 over the previous row's hash and the commitment. A purge (\lean{remove} of a path,
\lean{remove_digest} of a digest) empties a row: the path, the identity, the digest, the title,
the salt and the bookmark go, and the position, the kind, the key, the closing stamp, the
commitment and the hash stay. The chain then verifies from its first row as before, and the
commitment, its salt gone, tests no guess. The row keeps its path's key because a purge is by
digest and may take an old version of a path that is still live: a hole that lost its key would
let a read by path return the version before it at a cut where the purged one was the head. The
key is $\mathrm{HMAC}(k_c, \text{path})$ with $k_c = \mathrm{HMAC}(K, \text{chain's bytes})$,
where $K$ is the instance's landing-log secret, read from the environment as the store reads its
other secrets and never written to the file. The design had the chain's bytes alone as the key,
in the same file as the rows; with $K$ outside it, a copy of the file cannot test whether a hole
was a given path. A missing secret refuses the migration and every write before anything is
written, and a chain's key check refuses a secret it was not written with.

\begin{algorithm}[t]
\caption{A landing of \texttt{path} in class \texttt{cls}, inside the catalog's transaction.}
\label{alg:record}
\begin{algorithmic}[1]
\State $p \gets$ the catalog's row of \texttt{path}, read before it is rewritten
\State upsert the catalog's row (the key, as before)
\If{$p$ is active with the same digest, title, class and identity}
  \State \Return \Comment{a write that changes nothing appends nothing}
\EndIf
\State $e \gets$ the path's last event in the one chain where it is not a \texttt{leave}
\If{$e$ exists and its chain holds \texttt{cls}}
  \State $T \gets e$'s chain; \textbf{if} $e$ is a \texttt{store} of another identity \textbf{then}
    append a \texttt{move} to $T$
\Else
  \State $T \gets$ a chain of \texttt{cls} the path left, else the newest chain of \texttt{cls},
    else a new chain (\texttt{genesis}, \texttt{horizon})
  \State \textbf{if} $e$ exists \textbf{then} append a \texttt{leave} to $e$'s chain
\EndIf
\State append a \texttt{store} to $T$: set the path's open row's closing stamp to the new
  position, then insert the row with its salt, commitment and hash
\end{algorithmic}
\end{algorithm}

\paragraph{The genesis.}
On the first open of a catalog that has rows and no log, each class gets a chain: a
\lean{genesis}, the catalog's rows of the class in path order, an active path as a \lean{store}
and a tombstoned one as a \lean{store} followed by its tombstone, then a \lean{horizon} whose
position is its own. The order inside the block is not history, so the positions under the horizon
are not served at a cut. The design wrote every such tombstone as a \lean{retire}; the build writes
the \lean{move} that \lean{tombstone} would, where another active path holds the identity, so that
a renamed path is not counted among the identities whose history runs across a retire.

\paragraph{Reads.}
\Cref{alg:record} is the write. The reads are \lean{chains(class)}, \lean{head(chain)},
\lean{read_at} and \lean{read_between} (the rows of a range of positions), \lean{state_at} (a
path's last event at or under a cut over its class's chains, one seek on the key index a chain),
\lean{history_at} (its \lean{store} rows there, a hole where a version was purged),
\lean{heads_at} (the \lean{store} rows whose closing stamp is absent or past the cut), and
\lean{verify}, which recomputes the hashes from the first row and the commitment of every row that
still has its salt. A cut over a class is one position a chain. Nothing reads the log yet; the
stage that reads the notes at a cut does.

\section{Evaluation}\label{sec:eval}

\paragraph{What is proved.}
The 47 declarations the design lists for SNP-1 to SNP-5 are restated, their proofs left out, in a
statement file over the model, and a comparator checks that the library proves each exactly as
stated, from the standard axioms; the store's gate also audits the axioms of all 5{,}250 library
declarations and replays the export, 56{,}231 declarations from \lean{Init} up, through Lean's
kernel and five independent checkers. The ledger now has 105 rows in 16 families, and a
cross-reference holds its rows, the statements and the docstring markers to one another. The
attacks on these laws, property tests of the model's computable definitions run before the proofs,
survive 500 draws each.

\paragraph{What is tested.}
The transcription is checked against the laws by a property test: Hypothesis draws sequences of 3 to
40 catalog writes over five paths, six digests, two titles, three classes and two identities, with
retires, renames, relabels, re-keys and both purges among them, 200 sequences a run. After every
write the catalog must equal the fold of the log and every chain must verify; every cut taken during
the sequence must read, after every later write, the rows and the state of every path it read
then, up to the holes a later purge leaves (SNP-1); every read of a chain must be a prefix of the
read at its head (SNP-2); the fold of every cut taken, continued over the rows since, must be the
catalog (SNP-3); the state of a path at the head must be its catalog row (SNP-4); and at every
position of every chain the seek, the closing-stamp read and the fold of the prefix must give every
path the same row, with no path alive twice (SNP-5). Eleven more tests take the stage's cases: a
write that changes nothing, and a landing the writer's unchanged check skips, append nothing; a
process killed between the catalog's row and the log's leaves neither; a failed batch leaves no
log row; the genesis is the catalog, tombstoned paths included, and ends at a horizon; a relabel
into a class that has rows changes no row; a purge of an old version of a live path is a hole by
all three reads at a cut where it was the head; an edited row fails \lean{verify}; the secret is
required, and another is refused; and every statement by path is planned on the key index. Seeding
two bugs, a log that writes no closing stamp and one that writes no \lean{leave}, fails the suite.

\paragraph{What is measured.}
\Cref{tab:measure} gives two runs of the committed harness on two copy-on-write clones of the
catalog of the store's corpus instance, 400{,}274 rows of one class, all active, on an Apple M2
Pro with SQLite 3.54.1 and other work running; the original was never opened. The genesis is one
chain of 400{,}276 rows, and its fold equals the catalog.

\begin{table}[t]
\centering
\small
\begin{tabular}{@{}p{0.50\linewidth}p{0.44\linewidth}@{}}
\toprule
measurement & result \\
\midrule
the migration: tables, indexes, genesis & 6.22 to 6.33 s \\
bytes a row of the log (table, key index) & 429 (381, 47); the file from 162.7 to 334.1 MB \\
\lean{verify}, every commitment and hash & 3.05 to 3.07 s \\
\lean{state_at} of one path; its keyed hash & 14.3 to 14.8 $\mu$s; 1.6 $\mu$s \\
\lean{history_at} of one path & 14.7 to 15.7 $\mu$s \\
\lean{read_at}, \lean{read_between}, 1{,}000 positions & 3.3 to 3.4 ms \\
\lean{heads_at} at a tenth, the middle, the head & 0.15, 0.79 to 0.80, 1.6 s \\
the same walks counted in SQL & 6, 26, 54 to 57 ms \\
a new path's write, with the log and without & 49.0 and 6.8 $\mu$s \\
a new version's write, with the log and without & 106.0 to 106.5 and 37.9 to 41.2 $\mu$s \\
\bottomrule
\end{tabular}
\caption{The landing log on a copy of a 400{,}274-item catalog: two runs, medians, ranges across
the runs. Writes are in batches of 10{,}000, one transaction each, as a landing batch writes its
keys.}\label{tab:measure}
\end{table}

The log doubles the catalog's file, with paths, digests and titles the catalog holds already. It
adds 42~$\mu$s to a new path's write and 65 to 69~$\mu$s to a new version's, which also writes a
closing stamp; a landing without its vector took 31~ms an item on the store's earlier measurement
of two-phase ingestion, so the log is 0.1 to 0.2 percent of a landing. Reading every head at the
head of this chain takes 1.6~s, almost all of it building 400{,}274 rows in Python; the walk alone
is 0.14~$\mu$s an event. That is past the tenth of a second the design allows a read from the log
alone, and the design already keeps heads at recorded cuts for a chain over 250{,}000 events; this
chain is one.

\paragraph{What measuring found.}
The first run took about 40~ms a landing. The update that writes the closing stamp of a path's
previous row took the primary key, \lean{(chain, seq)}, and walked the chain for the path's rows,
though the statement names the path's key: the planner's choice the design had already met for a
seek on its fixture. The statement now names the key index, and a test asks the plan of every
statement by path to name it. Nothing at the test suite's sizes showed the walk.

\section{Limitations}

One writer is assumed, as by the model. A build that wrote the catalog without the log, while the
log's tables stayed, would leave the log behind the catalog, and nothing detects that yet; the
rollback drops the tables, and the next open writes a new genesis. The genesis knows no version
numbers, since the version chains are another file. A purge of a digest reads the whole log, which
has no index by digest. A path with events in two chains of one class, possible only after a
change of class and a merge of the two classes, is read by a rule of the prefix alone that can
order its earlier segments otherwise than they happened. The secret is an operational dependency:
a process that writes the catalog stops without it, and a lost secret leaves the log unwritable.
The measurements are of one instance, one class and no deletions; a catalog of several classes,
with identities apart from paths, is tested and not measured.

\section{Future work}

The next stages of the design number the chat log and the attestation ledger, record cuts and
resolve names to them; then take the two forms of removal into every store (the deep form's reach
and the fresh identity), then read the notes, vectors and lexical index at a cut from this log.
Each property-tests its transcription against the Lean statements, and its laws then enter the
ledger.

\section{Conclusion}

A catalog rewritten in place loses its past at every write. Two tables in its own file, written in
the same transactions, keep it at about 430 bytes and 50~$\mu$s a landing, with the catalog their
fold. Holes keep their positions and their keys, so the chain survives an erasure and a read by
path never steps over one; the key's other half stays outside the file, so the holes reveal
nothing to whoever copies it. The five laws a reader at a cut relies on are now statements the
library is judged against, and the code is tested against them.

\bibliographystyle{plainnat}
\bibliography{refs}
\end{document}
